\chapter{Pendahuluan}

\section{Latar Belakang}

Marisa sedang berusaha menyusun sebuah arsip rahasia yang memuat pemetaan
wilayah, entitas penguasa, dan kronologi insiden penting di Gensokyo. Berkas
ini merupakan gabungan dari catatan riset pribadinya serta dokumen-dokumen
sensitif yang berhasil ia kumpulkan dari berbagai pihak. Demi menjaga keamanan
informasi, seluruh naskah di dalam arsip tersebut dienkripsi dengan berbagai
jenis \textit{cipher} klasik. Sebagai rekan riset yang dipercaya, tugas kami
adalah membantu Marisa mendekripsi serangkaian dokumen dalam arsip ini agar
seluruh pengetahuan rahasia tersebut dapat terungkap.

\section{Ringkasan Soal}

Laporan ini terdiri atas lima soal wajib dan satu soal bonus. Setiap soal
menuntut kriptanalisis terhadap sebuah \textit{cipher} klasik yang berbeda.

\begin{table}[H]
\centering
\begin{tabular}{c l l}
\toprule
\textbf{No.} & \textbf{Judul} & \textbf{Algoritma \textit{Cipher}} \\
\midrule
2.1 & Beast Metropolis          & Substitusi abjad-tunggal (\textit{monoalphabetic}) \\
2.2 & Horizon Chartreuse        & Vigen\`ere \textit{Cipher} \\
2.3 & Domain of Hopes and Dreams & Playfair \textit{Cipher} \\
2.4 & Trinitarian Rhapsody      & Affine \textit{Cipher} \\
2.5 & Silent Goddess            & Hill \textit{Cipher} \\
\midrule
3.1 & Ex Machina (Bonus)        & Enigma Machine M3 (termodifikasi) \\
\bottomrule
\end{tabular}
\caption{Daftar soal dan algoritma \textit{cipher} yang digunakan.}
\end{table}

\section{Ketentuan Pengerjaan}

Sesuai spesifikasi tugas, seluruh proses enkripsi pada \textit{plaintext} asli
dilakukan setelah teks diubah menjadi huruf kapital dan hanya pada karakter
alfabetik \texttt{[A--Z]}; karakter lain (angka, spasi, tanda baca) dibuang dan
tidak dienkripsi. Untuk soal wajib, kriptanalisis dilakukan hanya dengan
\textit{approach} yang telah diajarkan di kelas. Bantuan skrip program, lembar
kerja, atau alat bantu komputasi lain diperbolehkan untuk menghitung frekuensi
karakter dan pola.
